
\begingroup
\newcommand{\SGAIIIIntroFootnote}[2]{%
  \begingroup
  \renewcommand{\thefootnote}{#1}%
  \footnote{#2}%
  \endgroup
}

\chapter*{Editorial Notice}
\addcontentsline{toc}{chapter}{Editorial Notice}
\markboth{EDITORIAL NOTICE}{EDITORIAL NOTICE}

We present here a slightly revised edition of the original Seminar, whose purpose
and contents are described in the Introduction.  The revision consisted mainly of
correcting typographical errors, adding a few supplementary remarks or references
in footnotes, dividing the work in its present form into three volumes, each with a
detailed table of contents and an index of notation, and adding a terminological index at
the end of Volume~3.  In addition, Expos\'e~VI\textsubscript{B}, by J.-E. Bertin,
was partially rewritten by the author, especially \SGARefLink{sga3:intro:section:5}{Sections~5} and~10, so some
references to that Expos\'e differ from the references to the original version.  A
list of the Expos\'es in the Seminar appears at the beginning of this volume.

Since the first edition of the present Seminar appeared, all of Chapter~IV of the
\emph{\'El\'ements de G\'eom\'etrie Alg\'ebrique} has been published, making certain
passages of the Seminar unnecessary.  We have occasionally indicated in footnotes
the relevant references to EGA~IV that allow those passages to be bypassed.

For another treatment of algebraic groups that systematically uses the language of
schemes, we refer to M. Demazure and P. Gabriel, \emph{Groupes Alg\'ebriques}
(North-Holland \& Masson et Cie).  Unlike the present Seminar, that book assumes no
knowledge of algebraic geometry, but includes all the required preliminaries from
scheme theory; it can therefore serve as an introduction to the study of this
Seminar.  (It also treats subjects not covered in the Seminar, such as the
Dieudonn\'e-style structure theory of commutative affine algebraic groups, in
\emph{loc. cit.}, Chapter~V.)

\begin{flushright}
Bures-sur-Yvette, March 1970
\end{flushright}

\clearpage
\thispagestyle{empty}
\mbox{}
\clearpage

\chapter*{Introduction}
\addcontentsline{toc}{chapter}{Introduction}
\markboth{INTRODUCTION}{INTRODUCTION}

\section*{\SGAInlineRefTarget{sga3:intro:section:1}1. The purpose of this Seminar is twofold.}

On the one hand, we aim to provide convenient foundations for the theory of group
schemes in general.  Expos\'es~I--IV will supply the indispensable preliminary
exercises in schematic and categorical syntax.  To obtain a language that fits
geometric intuition without effort, and to avoid circumlocutions that become
intolerable in the long run, we always identify a prescheme~$X$ over another
prescheme~$S$ with the functor
\[
  (\mathbf{Sch})_{/S}^{\circ}\longrightarrow \mathbf{Set}
\]
that it represents\SGAIIIIntroFootnote{\ensuremath{\ast}}{Such a viewpoint seems
to have been considered for the first time eight or nine years ago, in connection
with the theory of formal groups, by P. Cartier, who unfortunately did not take the
trouble to state it precisely and systematize it as it deserved.}.  It is therefore
necessary to formulate many definitions so that they apply to arbitrary functors,
whether representable or not.  Moreover, nearly all the functors we shall use will
be ``sheaves'' (for the ``faithfully flat quasi-compact topology'').  Expos\'e~IV,
which treats groups only incidentally, sketches the language of ``localization'' and
sheaves; this language also proves very convenient in questions concerning the
representability of functors.  Above all, that Expos\'e will provide the most
convenient framework for the questions of passage to a quotient that arise later.

Expos\'e~V gives some general results on the existence of quotients.  They are taken
up again in Expos\'e~VI\textsubscript{A} for the quotient of an algebraic group over
a field (or, more generally, over an Artinian ring) by a
subgroup\SGAIIIIntroFootnote{\ensuremath{\ast\ast}}{For a more thorough study of
passage to the quotient, especially by reductive groups, see D. Mumford's important
work \emph{Geometric Invariant Theory}, Ergebnisse der Mathematik, vol.~34,
Springer, 1965.  We note that on one important point the terminology of that book
does not agree with ours: over a field of characteristic $p>0$, the groups Mumford
calls ``reductive'' are the smooth groups of multiplicative
type in the sense of the Seminar (cf. Expos\'e~IX).  One may reasonably suppose
that Mumford adopted this meaning of ``reductive,'' which loses its significance
over a base that is not a field, provisionally and as a kind of stopgap (and this is
approximately what Mumford explains, for other reasons, as early as the second
paragraph of his preface).}.  That Expos\'e and the following
Expos\'e~VI\textsubscript{B} also contain various elementary results specific to
algebraic groups over a field that are used routinely later.

Expos\'e~VII studies certain phenomena related to the characteristic of the base
field and, in particular, develops with the appropriate generality the
correspondence between radicial group schemes of height~1 and restricted
$p$-Lie algebras.

Finally, Expos\'e~XVIII contains the scheme-theoretic generalization of Weil's
theorem on the ``birational'' definition of algebraic groups.

On the other hand, we intend to generalize the Borel--Chevalley structure theory of
affine algebraic groups to groups over an arbitrary base prescheme.  In fact, while
the Seminar notes were being written, it became apparent that the affine hypothesis
was unnecessary for many results of the theory.  The most complete results are of
course obtained for ``semisimple group schemes,'' or more generally ``reductive
group schemes,'' with which we shall be exclusively concerned beginning with
Expos\'e~XIX.  Chevalley himself had already constructed the ``T\^ohoku groups''
over the ring of integers, and that construction will be taken up in the present
Seminar.  The principal uniqueness
theorem gives a simple characterization of the ``twisted''
variants of these T\^ohoku groups over a base prescheme~$S$: they are the affine,
smooth group schemes over~$S$ whose geometric fibers are connected semisimple
groups in the usual sense\SGAIIIIntroFootnote{\ensuremath{\ast\ast\ast}}{This is
the essential result of M. Demazure's thesis, \emph{Sch\'emas en groupes
r\'eductifs}, Bull. Soc. Math. France \textbf{93} (1965), 369--413.}%
.

\section*{\SGAInlineRefTarget{sga3:intro:section:2}2.}

As in the familiar theory over an algebraically closed field, the subtori of the
group schemes under consideration play a crucial role.  Thus the preliminary study
of tori and, more generally, of ``group schemes of multiplicative type''---both
intrinsically and as multiplicative-type subgroups of a given group---occupies a
fairly large part of this Seminar (Expos\'es~VIII--XII).  Their remarkable rigidity
(in some respects even greater than that of abelian schemes or semisimple group
schemes) makes them highly effective tools for studying certain more general groups.

\section*{\SGAInlineRefTarget{sga3:intro:section:3}3.}

Beginning with Expos\'e~XII (apart from Expos\'e~XVIII, already mentioned), we shall
routinely use the theory of affine algebraic groups over an algebraically closed
field, which the reader will find in the 1956 Chevalley Seminar, especially in
Expos\'es~IV--IX of that Seminar.  We shall also use, though to a lesser extent, the
later Expos\'es of the Chevalley Seminar devoted to the structure of semisimple
algebraic groups.  Indeed, we shall redevelop Chevalley's theory directly within the
framework of schemes; it will be seen that this approach makes the exposition
simpler and more precise, even over a base field.

\section*{\SGAInlineRefTarget{sga3:intro:section:4}4.}

The principal purpose of this Seminar is, of course, to develop techniques that
apply to the study of group schemes over an arbitrary base, that is, essentially,
to the study of families of algebraic groups.  The infinitesimal properties of such
families, and in particular the case of an Artinian base scheme, therefore play an
important role.  These properties arise even in the study of algebraic groups over
a field~$k$ when the field is not perfect, especially when applying descent in the
non-Galois case.  Among the new results obtained in this setting, let us mention the
existence of maximal tori and Cartan subgroups in an arbitrary smooth algebraic
group, the rationality of the variety of maximal tori, and various related results
(Expos\'e~XIV); another example is the correspondence between the ``forms''
of a semisimple group and homogeneous principal bundles under a suitable semisimple
algebraic group (generally not connected) (\SGARefLink{sga3:xxiv:1.16}{Expos\'e~XXIV, 1.16}--\SGARefLink{sga3:xxiv:1.20}{1.20}).  In general,
one may say that the methods required for work over a nonperfect base field are
essentially those used over arbitrary base preschemes and thus lie outside the
framework of classical algebraic geometry.

\section*{\SGAInlineRefTarget{sga3:intro:section:5}5.}

It did not seem useful to indicate, at the beginning of the written Expos\'es, the
date or dates of the corresponding oral presentations in the Seminar.  We shall
merely say that the order of the mimeographed Expos\'es (I--XXVI) does indeed agree
with the order of the oral presentations.  Moreover, the final written text was
sometimes prepared considerably later than the oral presentation and often differs
from it quite substantially: the written text is generally more detailed and more
complete (as in Expos\'es~IV and VII\textsubscript{B}), or even appreciably more
general (as in Expos\'es~XII and VII\textsubscript{B}), than the oral presentation.
Other written Expos\'es correspond to no oral presentation
(VI\textsubscript{B}, VII\textsubscript{A}, XV, XVI, XVII, and most of XXVI); they
were written and inserted into the mimeographed Seminar either to provide convenient
references for various other Expos\'es (notably in the case of
VI\textsubscript{B}) or because they form a natural continuation of the concepts
and techniques already developed.  Consequently, Expos\'es~VII\textsubscript{A},
VII\textsubscript{B}, XV, XVI, and XVII need not be read in order to study the rest
of the Seminar, particularly the part devoted to reductive group schemes.

\section*{\SGAInlineRefTarget{sga3:intro:section:6}6.}

From scheme theory we shall use primarily the general language of schemes set out in
EGA~I; the notions of flat, \(\acute{e}\)tale, and smooth morphisms set out in
SGA~1, I--V; and the theory of faithfully flat descent in SGA~1, VIII.  As far as
possible, we have avoided imposing unnecessary Noetherian hypotheses.  In return,
this has required us to replace the usual hypothesis ``of finite type'' with the
hypothesis ``of finite presentation.''  For the notion of a morphism of finite
presentation, the reader should consult EGA~IV, 1.4 and~1.6.  The results of
SGA~1, I--IV, which are most often stated in the Noetherian setting, will be
developed in general in EGA~IV\SGAIIIIntroFootnote{\ensuremath{\ast\ast\ast\ast}}{Since
this Introduction was written, all four parts (Sections~1--21) of EGA~IV have
appeared.}; standard methods for reducing certain kinds of statements involving
finite-presentation hypotheses to the Noetherian case will also be developed there
in detail (EGA~IV, Sections~8, 9, and~11).  A reader who does not wish to assume
these results from EGA~IV may simplify certain statements or their proofs by
assuming the base prescheme to be locally Noetherian.  This does, however, lead to
difficulties when a proof proceeds by descent from~$\widehat A$ to~$A$, where
$\widehat A$ is the completion of a Noetherian local ring~$A$, because that method
introduces the generally non-Noetherian ring
\[
  \widehat A\otimes_A\widehat A.
\]

\section*{\SGAInlineRefTarget{sga3:intro:section:7}7.}

References follow the usual decimal system: the reference 5.7.11 denotes the
proposition (or lemma, definition, and so forth) with that number in the same
Expos\'e; in a reference such as XVII~7.8, the Roman numeral denotes the Expos\'e.
We use the following abbreviations for our standard references:

\begin{description}
\item[\emph{Bible}] Chevalley Seminar, \emph{Groupes de Lie alg\'ebriques},
  1956/58.
\item[EGA $X$, $x.y.z$] J. Dieudonn\'e and A. Grothendieck,
  \emph{\'El\'ements de G\'eom\'etrie Alg\'ebrique}, Chapter~$X$, statement
  $x.y.z$ (or subsection $x.y$, and so forth).
\item[SGA $n$, $X$ $y.z$] \emph{S\'eminaire de G\'eom\'etrie Alg\'ebrique du
  Bois-Marie}, year~$n$, statement $y.z$ of Expos\'e~$X$.
\item[TDTE] A. Grothendieck, \emph{Techniques de descente et Th\'eor\`emes
  d'existence en G\'eom\'etrie Alg\'ebrique}, Expos\'es delivered in the Bourbaki
  Seminar between 1959 and 1962.
\end{description}

\endgroup
